% \section{Theoretical Analysis}
% \label{sec:theory}

% This appendix provides formal analysis of the V1 LP formulation.

% \subsection{Exactness of McCormick Linearization}

% \begin{proposition}
% \label{prop:mccormick_exact}
% Let $y \in \{0,1\}$ and $p \in \{0,1\}$. Define $w$ by the McCormick constraints:
% $w \leq y$, $w \leq p$, $w \geq y + p - 1$, $w \geq 0$.
% Then $w = y \cdot p$.
% \end{proposition}

% \begin{proof}
% There are four cases. If $y=0$: the constraints give $w \leq 0$ and $w \geq 0$, so $w = 0 = 0 \cdot p$. If $p=0$: symmetrically, $w = 0$. If $y = p = 1$: $w \leq 1$, $w \geq 1$, so $w = 1$. This covers all cases.
% \end{proof}

% This means that at any integer-feasible point of the V1 LP, the McCormick auxiliary variables $w_{i,c,k}$ exactly represent the bilinear products $y_{i,c} \cdot p_{\text{assoc}(i,k), o(i)}$, and the LP objective value equals the true joint tiling-fusion objective.

% \subsection{LP Relaxation Quality}

% When solved as a continuous LP (all variables in $[0,1]$), the McCormick linearization provides the tightest possible convex relaxation of the bilinear term over the unit box.

% \begin{proposition}
% The McCormick envelope of $f(y,p) = y \cdot p$ over $[0,1]^2$ is the convex envelope.
% \end{proposition}

% \begin{proof}
% The convex envelope of a bilinear function over a box is the maximum of the two underestimators $\max(0, y + p - 1)$, which is exactly the McCormick lower bound. Similarly, the concave envelope is $\min(y, p)$. This is a classical result (McCormick, 1976).
% \end{proof}

% This implies that the V1 LP relaxation provides the best possible lower bound achievable by any convex relaxation of the bilinear tiling-fusion interaction. The relaxation gap (difference between LP optimum and integer optimum) depends on problem structure; we report empirical gap measurements in the experiments.

% \subsection{Structural Invariance of the Constraint Matrix}

% \begin{proposition}
% \label{prop:invariance}
% For a fixed neural network model with $L$ operations, $T$ timesteps, $|\mathcal{E}|$ edges, and $|\mathcal{C}_i|$ Pareto configurations per operation, the sparsity pattern (row indices, column indices) of the V1 constraint matrix $A$ is invariant across all hardware configurations proposed by Vizier. Only the numerical values of $A$, the objective vector $c$, and the constraint bounds $l_c, u_c$ change.
% \end{proposition}

% \begin{proof}
% The constraint matrix structure is determined by the graph topology (which defines edges $\mathcal{E}$, execution order $o(i)$, and fan-in/fan-out sets), the number of Pareto configurations $|\mathcal{C}_i|$, and the timestep count $T$---all of which depend only on the neural network model, not on hardware parameters. Hardware parameters affect $\Tmax_i(c)$, $\Tmin_i(c)$, $t_i^k(c)$, $d_e$, $d_i^W$, $\CGM$, and $c_i^{\text{remat}}$, which enter as numerical coefficients. The assignment, McCormick, capacity, persistence, and boundary constraints all have the same variable participation pattern regardless of these coefficients.
% \end{proof}

% This invariance has practical implications: LP solver preconditioners can be factored once per model and reused across all hardware trials, and warm-starting from a previous trial's solution is straightforward since the variable ordering is stable.

% \subsection{NNZ Breakdown}

% Table~\ref{tab:nnz_breakdown} provides a detailed breakdown of nonzeros in the V1 constraint matrix.

% \begin{table}[h]
% \centering
% \caption{NNZ breakdown by constraint block for V1 on EfficientNet-B7 ($L=400$, $T=400$, $|\mathcal{C}_i|=1{,}000$).}
% \label{tab:nnz_breakdown}
% \small
% \begin{tabular}{lrl}
% \toprule
% \textbf{Constraint Block} & \textbf{Est. NNZ} & \textbf{Source} \\
% \midrule
% One-config-per-op & 400,000 & $L \times |\mathcal{C}|$ \\
% Exec.\ time upper bound & 4,800,000 & $|\mathcal{C}| \times 4$ terms per op \\
% Exec.\ time lower bound & 400,000 & $|\mathcal{C}|$ terms per op \\
% McCormick linearization & 7,200,000 & $3 \times 2 \times L|\mathcal{C}| \times 3$ \\
% Global memory capacity & 320,000 & $T \times (|\mathcal{E}| + L)$ \\
% Activation persistence & 800,000 & $|\mathcal{E}| \times T \times 2$ \\
% Weight monotonicity & 320,000 & $L \times T \times 2$ \\
% Remat feasibility & 640,000 & Several $L \times T$ blocks \\
% Boundary conditions & 800,000 & $|\mathcal{E}| \times T + L \times T$ \\
% \midrule
% \textbf{Total} & $\mathbf{\sim 15.7\text{M}}$ & \\
% \bottomrule
% \end{tabular}
% \end{table}

% \subsection{Complexity of the Pareto Pre-Pass}

% \begin{proposition}
% The precomputation cost of the Pareto pruning pass is $O(L \cdot |\mathcal{C}_{\text{raw}}| \cdot t_{\text{Timeloop}})$ per hardware configuration class, where $|\mathcal{C}_{\text{raw}}| = O(10^7)$ is the raw configuration count and $t_{\text{Timeloop}}$ is the average Timeloop evaluation time. With bandwidth-invariant Pareto fronts, this cost is amortized across all Vizier trials sharing the same compute configuration.
% \end{proposition}

% For EfficientNet-B7: $400 \times 1{,}000 = 400{,}000$ Timeloop evaluations (after pruning), at ${\sim}5$ seconds each, totaling ${\sim}23$ hours on a single machine or ${\sim}30$ minutes with 512-way parallelism.

% page edit

\section{Theoretical Analysis}
\label{sec:theory}

This appendix provides formal analyses of the CHEETAH-V1 LP formulation.

\subsection{Exactness of McCormick Linearization}

\begin{proposition}
\label{prop:mccormick_exact}
Let $y \in \{0,1\}$ and $p \in \{0,1\}$. Define $w$ by the McCormick constraints:
$w \leq y$, $w \leq p$, $w \geq y + p - 1$, $w \geq 0$.
Then $w = y \cdot p$.
\end{proposition}

\begin{proof}
There are four cases. If $y=0$: the constraints give $w \leq 0$ and $w \geq 0$, so $w = 0 = 0 \cdot p$. If $p=0$: symmetrically, $w = 0$. If $y = p = 1$: $w \leq 1$, $w \geq 1$, so $w = 1$. This covers all cases.
\end{proof}

This means that at any integer-feasible point of the V1 LP, the McCormick auxiliary variables $w_{i,c,k}$ exactly represent the bilinear products $y_{i,c} \cdot p_{\text{assoc}(i,k), o(i)}$, and the LP objective value equals the true joint tiling-fusion objective.

\subsection{LP Relaxation Quality}

When solved as a continuous LP (all variables in $[0,1]$), the McCormick linearization is the tightest convex relaxation of \emph{an individual} bilinear term over the unit box.

\begin{proposition}
The McCormick envelope of $f(y,p) = y \cdot p$ over $[0,1]^2$ is the convex envelope.
\end{proposition}

\begin{proof}
The convex envelope of a bilinear function over a box is the maximum of the two underestimators $\max(0, y + p - 1)$, which is exactly the McCormick lower bound. Similarly, the concave envelope is $\min(y, p)$. This is a classical result~\citep{mccormick1976computability}.
\end{proof}

\paragraph{This does not extend to the joint problem, and we do not claim that it does.}
It is tempting to read the proposition as saying that the V1 relaxation is the
best possible convex relaxation of the tiling-fusion interaction. It does not,
and the distinction is worth stating precisely, because the envelope is a
statement about one product in isolation and the program contains a
\emph{family} of products sharing a one-hot. Take two configurations with
$y_1+y_2=1$ and $w_j=y_j p$. Every integral point satisfies $w_1+w_2=p$. Yet
%
\begin{equation}
\label{eq:rlt_counterexample}
y_1=y_2=p=\tfrac12,\qquad w_1=w_2=\tfrac12
\end{equation}
%
satisfies all four McCormick inequalities for each $j$ while giving
$w_1+w_2=1\neq p$. The per-term envelopes therefore admit points that no
integral solution can reach, and additional coupling constraints can strengthen
the relaxation \emph{without adding variables}. Since $\sum_c y_{i,c}=1$ already
holds, the reduction-and-linearization (RLT) equalities
%
\begin{equation}
\label{eq:rlt_linking}
\sum_c w_{i,c,k} = p_{\mathrm{assoc}(i,k)},\qquad
\sum_c z^{\mathrm{act}}_{e,c,t} = p_{e,t},\qquad
\sum_c z^{W}_{i,c,t} = p^{W}_{i,t}
\end{equation}
%
hold at every integral point and are valid cuts~\citep{sherali1990hierarchy}. We
include them. The correct statement of the guarantee is therefore: the V1
relaxation is exact at every integral vertex (Proposition~\ref{prop:mccormick_exact}),
and is the convex hull of each product \emph{individually}, strengthened by
\cref{eq:rlt_linking}; it is not the convex hull of the joint feasible set, and
no claim in this paper depends on it being so.

\paragraph{How much looseness this costs, measured.}
Because \cref{eq:rlt_linking} are valid cuts, adding them can only raise the LP
optimum, so the question is quantitative. We measured it by solving each program
twice, with and without the linking rows, holding everything else fixed. On the
V1 tile-dependent LP the two objectives agree \emph{bit-for-bit} at every
on-chip-memory size we tested ($1$, $2$, $4$, $8$ and $16$~MiB on
EfficientNet-B7: $3.4001876777\times10^{10}$ in all ten solves, despite the cuts
adding $112{,}746$ rows), and likewise on the serving program at every batch
size of \cref{tab:batch}. The reason is structural rather than fortunate: the
rows these cuts tighten are the on-chip capacity rows, and on these workloads
those rows are slack by three to four orders of magnitude. EfficientNet-B7's
optimum sits at $1.0041\times$ its pure compute floor while its DRAM bound is
$1.1\times10^{-4}$ of that floor, and the serving program independently selects
the \emph{smallest} on-chip memory on its menu, which is what a program does
when capacity is not scarce. We therefore report zero measured relaxation
looseness from this source on our suite, and state plainly that this is a
property of the compute-bound regime we evaluate, not a general guarantee: on a
capacity-bound workload \cref{eq:rlt_counterexample} would bite, and
\cref{eq:rlt_linking} is what prevents it. The relaxation gap between LP optimum
and integer optimum is reported empirically per workload in the experiments.

\subsection{Structural Invariance of the Constraint Matrix}
\label{app:invar}

\begin{proposition}
\label{prop:invariance}
For a fixed neural network model with $L$ operations, $T$ timesteps, $|\mathcal{E}|$ edges, and $|\mathcal{C}_i|$ Pareto configurations per operation, the sparsity pattern (row indices, column indices) of the V1 constraint matrix $A$ is invariant across all hardware configurations proposed by Vizier. Only the numerical values of $A$, the objective vector $c$, and the constraint bounds $l_c, u_c$ change.
\end{proposition}

\begin{proof}
The constraint matrix structure is determined by the graph topology (which defines edges $\mathcal{E}$, execution order $o(i)$, and fan-in/fan-out sets), the number of Pareto configurations $|\mathcal{C}_i|$, and the timestep count $T$, all of which depend only on the neural network model, not on hardware parameters. Hardware parameters affect $\Tmax_i(c)$, $\Tmin_i(c)$, $t_i^k(c)$, $d_e$, $d_i^W$, $\CGM$, and $c_i^{\text{remat}}$, which enter as numerical coefficients. The assignment, McCormick, capacity, persistence, and boundary constraints all have the same variable participation pattern regardless of these coefficients.
\end{proof}

This invariance has practical implications: LP solver preconditioners can be factored once per model and reused across all hardware trials, and warm-starting from a previous trial's solution is straightforward since the variable ordering is stable.

\subsection{NNZ Breakdown}

Table~\ref{tab:nnz_breakdown} provides a detailed breakdown of nonzeros in the V1 constraint matrix.

\begin{table}[h]
\centering
\caption{NNZ breakdown by constraint block for V1 on EfficientNet-B7 ($L=400$, $T=400$, $|\mathcal{C}_i|=1{,}000$).}
\label{tab:nnz_breakdown}
\small
\begin{tabular}{lrl}
\toprule
\textbf{Constraint Block} & \textbf{Est. NNZ} & \textbf{Source} \\
\midrule
One-config-per-op & 400,000 & $L \times |\mathcal{C}|$ \\
Exec.\ time upper bound & 4,800,000 & $|\mathcal{C}| \times 4$ terms per op \\
Exec.\ time lower bound & 400,000 & $|\mathcal{C}|$ terms per op \\
McCormick linearization & 7,200,000 & $3 \times 2 \times L|\mathcal{C}| \times 3$ \\
Global memory capacity & 320,000 & $T \times (|\mathcal{E}| + L)$ \\
Activation persistence & 800,000 & $|\mathcal{E}| \times T \times 2$ \\
Weight monotonicity & 320,000 & $L \times T \times 2$ \\
Remat feasibility & 640,000 & Several $L \times T$ blocks \\
Boundary conditions & 800,000 & $|\mathcal{E}| \times T + L \times T$ \\
\midrule
\textbf{Total} & $\mathbf{\sim 15.7\text{M}}$ & \\
\bottomrule
\end{tabular}
\end{table}

\subsection{Complexity of the Pareto Pre-Pass}

\begin{proposition}
The precomputation cost of the Pareto pruning pass is $O(L \cdot |\mathcal{C}_{\text{raw}}| \cdot t_{\text{Timeloop}})$ per hardware configuration class, where $|\mathcal{C}_{\text{raw}}| = O(10^7)$ is the raw configuration count and $t_{\text{Timeloop}}$ is the average Timeloop evaluation time. With bandwidth-invariant Pareto fronts, this cost is amortized across all Vizier trials sharing the same compute configuration.
\end{proposition}

For EfficientNet-B7: $400 \times 1{,}000 = 400{,}000$ Timeloop evaluations (after pruning), at ${\sim}5$ seconds each, totaling ${\sim}23$ hours on a single machine or ${\sim}30$ minutes with 512-way parallelism.

\subsection{Containment of the Sequential Map--Fuse--Size Pipeline}
\label{app:containment}
We show that CHEETAH's joint feasible polytope \emph{contains} every in-scope
sequential map-then-fuse-then-size pipeline, lifted FAST included: under matched
hardware parameters and mapper constants, every point feasible for any such pipeline embeds
objective-preservingly into CHEETAH's joint LP feasible set. CHEETAH therefore
\emph{strictly generalizes} the sequential pipeline rather than merely extending
it. Formally, we prove the one-directional statement (\emph{for every
sequential-pipeline feasible point there exists a CHEETAH-feasible extension
with equal objective}) rather than a set equality; the extra CHEETAH
coordinates (a free tile selection, rematerialization, and the CGM one-hot)
reach points the sequential pipeline cannot express. The result is a modeling
statement about feasible sets, not a prediction that the two systems report the
same numerical speedups.

Consider a chip $H$ (MAC count, DRAM bandwidth $b$, $C_{GM}$ size $g$, datatype, batch) and a single workload DAG $G$ with $L$ layers, $E$ edges, $T$ scheduling steps, and constants $T_i^{\min}(c), T_i^{\max}(c), t_i^k(c)$ at $C$ Pareto tile configurations per layer. 

Both modeling formulations include decision variables:
\begin{itemize}\setlength{\itemsep}{0pt}\setlength{\topsep}{0pt}
  \item $y \in [0,1]^{L \times C}$ -- tile-config selection per layer, $\sum_c y_{i,c}{=}1$,
  \item $p \in [0,1]^{E \times T}$ -- activation residency per (edge, timestep),
  \item $p_W \in [0,1]^{L \times T}$ -- weight residency per (layer, timestep),
  \item $T \in \mathbb{R}_{\ge 0}^{L}$ -- per-layer cycle count (objective summand).
\end{itemize}
The V1 LP additionally inherits $r \in [0,1]^{L \times T}$ (rematerialization) and
$h_{\mathrm{GM}} \in [0,1]^{K}$ with $\sum_k h_{\mathrm{GM},k}{=}1$ (one-hot global memory selection from a
discrete grid $g_1,\ldots,g_K$). The primary cycle objective is $\mathbf{c}^\top x = \sum_i T_i$.

\paragraph{In-scope sequential pipelines.}
\begin{definition}[Sequential map-then-fuse pipeline]
\label{def:seq_scope}
Fix $H$, $G$, a topological execution order, and the Timeloop cost constants.
A \emph{sequential map-then-fuse pipeline} is any procedure that
(i) in a \emph{mapping stage} commits every layer $i$ to a single tile
configuration $c_i \in \mathcal{C}_i$ from the shared Pareto menu, by any
mapper (Timeloop, ZigZag, CoSA, or a hand schedule; a tile outside the menu is
admitted by adding its cost constants as a column), and only then
(ii) in a \emph{fusion stage} chooses a static residency pattern $(p, p_W)$
under the shared $C_{GM}$-capacity and bandwidth model, with the global-memory
size fixed at some $g \in \{g_1,\ldots,g_K\}$ and without rematerialization.
Writing $c = (c_1,\ldots,c_L)$, its lifted feasible set in the V1 variable
space is
\begin{equation}
\begin{split}
\mathcal{F}_S^{\mathrm{lift}}(H,G;c,g) \;=\; \Big\{(y, p, p_W, T) \;\Big|\;
&\; y_{i,\, c_i} = 1,\quad
y_{i,\, c'} = 0 \;\;\forall\, c' \ne c_i, \\
&\; \text{$(p,p_W)$ static, chosen by the fusion stage,}\quad
A_S\, x \le b_S(H,G;g) \Big\},
\end{split}
\end{equation}
where $A_S$, $b_S$ encode the $C_{GM}$ capacity (with $g$ hard-coded),
bandwidth, and per-stage compute floor under the fixed-$y$ projection.
\end{definition}
Out of scope, by construction, are methods that change the execution
\emph{order} (the polyhedral reordering axis, Section~\ref{sec:related}), use a
cost model other than the shared roofline/Timeloop constants, or size hardware
outside the menu. FAST is the instance $c = c^{\mathrm{FAST}}$ (Timeloop's
per-op choice) with $g$ its fixed $C_{GM}$:
$\mathcal{F}_F^{\mathrm{lift}}(H,G) = \mathcal{F}_S^{\mathrm{lift}}(H,G;c^{\mathrm{FAST}},g^{\mathrm{FAST}})$.
FAST commits to $r{=}0$ and to fixed $g$ by construction.

\paragraph{CHEETAH-V1 LP relaxation $\mathcal{P}_V(H,G)$.}
$$
\mathcal{P}_V(H,G) \;=\; \{(y, p, p_W, r, h_{\mathrm{GM}}, T) \in [0,1]^{LC+ET+2LT+K} \times \mathbb{R}_{\ge 0}^{L} \mid A_V\, x \le b_V(H,G)\}.
$$
The first $LC + ET + LT$ columns of $A_V$ coincide with $A_S$. The bandwidth
constraint uses the same $b$, and the $C_{GM}$-capacity row is parameterized by
the chosen $g_k$ via $h_{\mathrm{GM}}$.

\begin{theorem}[Containment of every in-scope sequential pipeline in $\mathcal{P}_V$]
\label{thm:containment}
Fix $H$, $G$, the execution order, and the Timeloop constants. For
\emph{every} tile commitment $c \in \prod_i \mathcal{C}_i$ and every
$g \in \{g_1,\ldots,g_K\}$ there exists a lifting map
$$
\iota:\; \mathcal{F}_S^{\mathrm{lift}}(H,G;c,g) \;\hookrightarrow\; \mathcal{P}_V(H,G),\qquad
\iota(y,p,p_W,T) \;=\; \big(y,\,\tilde{p},\,p_W,\,r{=}0,\,h_{\mathrm{GM}}{=}e_{k^\star},\,T\big),
$$
with $k^\star$ satisfying $g_{k^\star}{=}g$ and $\tilde{p}_{e,t}{=}1$ exactly on
the timesteps covered by the fusion stage's region for the edge supplying $e$,
zero elsewhere. The map is well-defined, injective, and preserves the primary
cycle objective: $\sum_i T_i \circ\iota(x) = \sum_i T_i(x)$ for every
$x \in \mathcal{F}_S^{\mathrm{lift}}(H,G;c,g)$. In particular, lifted FAST
embeds in $\mathcal{P}_V$.
\end{theorem}

\begin{proof}
The integer cube embeds into the LP cube, $\{0,1\}^{d} \subset [0,1]^{d}$,
therefore any integer-feasible $(y,p,p_W)$ lies in $\mathcal{P}_V$ along the
shared coordinates.
The lifted $\tilde{p}$ equals $1$ from the producer step through the last step
of the fused region and $0$ afterwards, hence it is monotone non-increasing
after production; V1's activation-persistence rows
$p_{e,t} \le p_{e,t-1} + r_{\mathrm{src}(e),t}$ therefore hold with $r{=}0$,
and the capacity rows charge exactly the tensors the fusion stage keeps
resident, so they coincide with $A_S x \le b_S$. Setting $r{=}0$ disables V1's
rematerialization slack. Setting $h_{\mathrm{GM}}{=}e_{k^\star}$ pins V1's
$C_{GM}$ constraint to the pipeline's fixed $g$.
Moreover, the per-stage compute floor evaluates identically because both
formulations consume the same cost vectors under the singleton-$y$ projection.
Therefore, $\iota(x) \in \mathcal{P}_V$ for all
$x \in \mathcal{F}_S^{\mathrm{lift}}(H,G;c,g)$. Nothing in the argument depends
on \emph{how} $c$ or $(p,p_W)$ were chosen, only on $c_i \in \mathcal{C}_i$, on
the pattern being static, and on $g$ lying in the menu; the statement therefore
holds uniformly over the whole class of Definition~\ref{def:seq_scope}.

\emph{Injectivity.} $\iota$ is the identity on the shared coordinates and
constant on the new coordinates, thus $\iota(x){=}\iota(x') \Rightarrow x{=}x'$.

\emph{Primary-objective preservation.} The primary cycle objective
$\sum_i T_i$ depends only on the shared coordinate $T$, thus
$\mathbf{c}^\top \iota(x) = \mathbf{c}^\top x$.

If V1's objective additionally includes an $\varepsilon$-tiebreaker term
$\varepsilon\sum_k k\, h_{\mathrm{GM},k}$, then on $\iota(x)$ this term equals
the constant $\varepsilon k^\star$, which does not affect the primary cycle
optimum or its $\arg\min$ for fixed $g$.
\end{proof}

\paragraph{Strict generalization, stated crisply.}
Theorem~\ref{thm:containment} reads: \emph{for every in-scope sequential
pipeline (Definition~\ref{def:seq_scope}), every feasible point it can produce
has a CHEETAH-feasible extension with equal primary objective.} CHEETAH's
feasible set is therefore a superset of the union over all such pipelines,
$$
\bigcup_{c,\,g}\;\iota\big(\mathcal{F}_S^{\mathrm{lift}}(H,G;c,g)\big)
\;\subseteq\; \mathcal{P}_V(H,G),
$$
which is the precise sense in which sequential map-then-fuse methods search a
subset of CHEETAH's decision space. This is an objective-preserving embedding,
\emph{not} a set equality: CHEETAH carries strictly more degrees of freedom (a
free tile selection $y$, rematerialization $r$, and the CGM one-hot
$h_{\mathrm{GM}}$), so each lifted pipeline point lands on a face-restricted
slice of $\mathcal{P}_V$ and CHEETAH additionally reaches points no sequential
pipeline can express. Every map-then-fuse-then-size schedule any in-scope
pipeline can produce is thus reachable inside CHEETAH's single joint program at
the same cost. Moreover, because every lifted point is integral in $y$, and the
refinements that distinguish the richer CHEETAH formulations either add
collapsible variables (V2's per-edge residency and CGM one-hot) or replace a
bound that is already exact at integer $y$ (the disaggregated epigraph of
\S\ref{sec:v1}, cf.\ Corollary~\ref{cor:v3_hull}), the same embedding $\iota$
places every in-scope pipeline inside every CHEETAH formulation V1.
Containment is therefore a property of the joint program, not of one particular
variant.

\begin{corollary}[CHEETAH-V1 LP lower bound dominates every in-scope sequential pipeline]
\label{cor:optimum}
For every chip $H$, workload $G$, matched Timeloop constants, every tile
commitment $c$, and every menu size $g$,
$$
\min_{x \in \mathcal{P}_V(H,G)}\!\!\mathbf{c}^\top x \;\;\le\;\;
\min_{x \in \mathcal{F}_S^{\mathrm{lift}}(H,G;c,g)}\!\!\mathbf{c}^\top x,
\qquad\text{in particular}\qquad
\min_{x \in \mathcal{P}_V(H,G)}\!\!\mathbf{c}^\top x \;\;\le\;\;
\min_{x \in \mathcal{F}_F^{\mathrm{lift}}(H,G)}\!\!\mathbf{c}^\top x.
$$
Any fixed-hardware performance gap between a sequential pipeline and CHEETAH
V1 therefore cannot be attributed to sequential solutions lying outside the V1
search space, and must be due to (i) V1's extra continuous degrees of freedom
(free $y$, $r$, $h_{\mathrm{GM}}$), (ii) LP relaxation versus integer rounding,
or (iii) differences in the cost model or baseline normalization.
The corollary bounds the V1 LP relaxation, not the rounded executable schedule
(which is reported separately).
\end{corollary}
\begin{proof}
Fix any $c$ and $g$ and let $x_S^\star$ attain the minimum over $\mathcal{F}_S^{\mathrm{lift}}(H,G;c,g)$. By Theorem~\ref{thm:containment},
$\iota(x_S^\star) \in \mathcal{P}_V$ and $\mathbf{c}^\top \iota(x_S^\star) = \mathbf{c}^\top x_S^\star$; the FAST case is $c = c^{\mathrm{FAST}}$.
The V1 LP minimum is at most this feasible value.
\end{proof}

The containment above bounds how CHEETAH relates to the sequential pipeline. The
next corollary bounds how CHEETAH's own formulations relate to one another,
tying the certified lower bound to the disaggregated perspective epigraph of
CHEETAH-V1 (\S\ref{sec:v1}, \eqref{eq:v3_epigraph}).

\begin{corollary}[V3's disaggregated epigraph is the exact hull and tightens the certified bound]
\label{cor:v3_hull}
Fix the hardware ($v,u$) and mapper constants. Let the aggregate roofline row of
CHEETAH-V1 \eqref{eq:v1} be
$T_i \ge \max\!\big(\sum_c \Tmin_i(c)\,v_{i,c},\; \sum_c (\Tmax_i(c){-}\mathrm{save}_{i,c})\,u_{i,c}\big)$,
and let \eqref{eq:v3_epigraph} be the disaggregated perspective form with split
variables $v_{i,c},u_{i,c},\tau_{i,c}$ obeying $\sum_c v_{i,c}{=}v$,
$\sum_c u_{i,c}{=}u$ and the box links
$\underline v\,y_{i,c}\le v_{i,c}\le \overline v\,y_{i,c}$,
$\underline u\,y_{i,c}\le u_{i,c}\le \overline u\,y_{i,c}$.
Then \eqref{eq:v3_epigraph} is the closed convex hull of the per-op disjunction
``select one tile $c$ and pay its roofline maximum
$\max(\Tmin_i(c)\,v,\,(\Tmax_i(c){-}\mathrm{save}_{i,c})\,u)$.'' Consequently:
\begin{enumerate}\setlength{\itemsep}{0pt}\setlength{\topsep}{0pt}
  \item[(i)] \eqref{eq:v3_epigraph} is \emph{exact at every integer vertex}
  $y\in\{0,1\}$, where it collapses to the true selected-tile roofline and hence
  coincides with \eqref{eq:v1};
  \item[(ii)] as a continuous relaxation it enforces
  $T_i \ge \sum_c \max(\Tmin_i(c)\,v_{i,c},\,(\Tmax_i(c){-}\mathrm{save}_{i,c})\,u_{i,c})$,
  which \emph{dominates} the weaker aggregate bound
  $T_i \ge \max(\sum_c \Tmin_i(c)\,v_{i,c},\,\sum_c (\Tmax_i(c){-}\mathrm{save}_{i,c})\,u_{i,c})$
  of \eqref{eq:v1}.
\end{enumerate}
Therefore, at matched hardware, the CHEETAH-V1 certified lower bound is never
below the CHEETAH-V1 lower bound,
$\min_{\mathcal{P}_{V3}}\mathbf{c}^\top x \ge \min_{\mathcal{P}_{V2}}\mathbf{c}^\top x$,
with strict improvement exactly when the V2 optimum blends a compute-bound tile
against a bandwidth-bound tile to hide both bottlenecks. The certified bound thus
\emph{cannot get worse and generally improves}, while both remain $\le$ the
lifted-FAST optimum of Corollary~\ref{cor:optimum}.
\end{corollary}
\begin{proof}
The box links are the McCormick envelope of the products $v\cdot y_{i,c}$ and
$u\cdot y_{i,c}$ over the unit-$y$ box, which is the tightest convex relaxation
and is exact at $y\in\{0,1\}$~\citep{mccormick1976computability}. At an integer
vertex, $\sum_c y_{i,c}=1$ forces $y_{i,c^\star}=1$ for a single $c^\star$ and
$y_{i,c}=0$ otherwise; the box links then pin $v_{i,c^\star}=v$, $u_{i,c^\star}=u$
and $v_{i,c}=u_{i,c}=0$ for $c\ne c^\star$, so
$\tau_{i,c^\star}\ge\max(\Tmin_i(c^\star)v,(\Tmax_i(c^\star){-}\mathrm{save}_{i,c^\star})u)$
and $T_i\ge\sum_c\tau_{i,c}$ reduces to the true selected-tile roofline, which is
\eqref{eq:v1} at that vertex, giving (i). For (ii), the epigraph rows
$\tau_{i,c}\ge\Tmin_i(c)\,v_{i,c}$ and
$\tau_{i,c}\ge(\Tmax_i(c){-}\mathrm{save}_{i,c})\,u_{i,c}$ give
$\tau_{i,c}\ge\max(\cdot,\cdot)$, so
$T_i\ge\sum_c\max(\cdot,\cdot)\ge\max(\sum_c\cdot,\sum_c\cdot)$ by superadditivity
of the maximum ($\sum_c\max(a_c,b_c)\ge\max(\sum_c a_c,\sum_c b_c)$). Thus every
$T$ feasible for \eqref{eq:v3_epigraph} satisfies the aggregate row of
\eqref{eq:v1}, i.e.\ the V3-feasible set is contained in the V2-feasible set
along the shared coordinates; minimizing the same objective over a subset yields
a value at least as large, giving
$\min_{\mathcal{P}_{V3}}\mathbf{c}^\top x \ge \min_{\mathcal{P}_{V2}}\mathbf{c}^\top x$.
By (i) the tightening removes only fractional-$y$ points, so lifted FAST, being
integral in $y$, stays feasible in $\mathcal{P}_{V3}$; hence both minima remain
$\le$ the lifted-FAST optimum via Corollary~\ref{cor:optimum}.
\end{proof}

Theorem~\ref{thm:containment} is a modeling containment result, and does not predict numerical equivalence of the speedups CHEETAH and FAST report.
Empirically, differences can arise from baseline normalization, TPU-v3 simulator overheads, control/padding costs, and batch/core conventions.

\begin{remark}[Scope of the certificate: inner LP vs.\ outer search]
\label{rem:scope}
The guarantees above are \emph{inner-problem} certificates. For a fixed hardware
candidate $H$, workload $G$, execution order, and finite candidate tiling sets,
CHEETAH's joint program is a convex LP, so any converged optimum is global to
within the reported solver tolerance and yields a certified lower bound on the
achievable schedule cost for that $H$. The \emph{outer} hardware search (the
LLM-Architect / Vizier loop that proposes chip candidates) is a heuristic: we do
\emph{not} claim a certified global optimum over the unrestricted hardware space,
only over the schedule for each proposed $H$.

The residual LP-to-integer gap observed on weight-heavy models, where a
per-block weight tensor exceeds the on-chip budget and cannot be held
integrally, is \emph{certified relaxation looseness}: a property of the
\emph{bound}, not of the deployed schedule. On the chain-structured subproblems
where it appears, the executable schedule attains the exact integral optimum
(verified by chain dynamic programming), so the gap measures the distance from
the LP lower bound up to that integral optimum, not any suboptimality of what
runs on hardware. In particular, the zero LP \emph{duality} gap at convergence
must not be conflated with this LP-to-integer \emph{relaxation} gap: the former
certifies the relaxation was solved to optimality, whereas the latter is the
inherent looseness of the convex relaxation, precisely the quantity that
Corollary~\ref{cor:v3_hull}'s tighter epigraph reduces.
\end{remark}